\documentclass[10pt,a4paper]{amsart}
\usepackage[margin=19mm]{geometry}
\usepackage{amsmath,amssymb,mathtools}
\usepackage[scr=rsfs,cal=euler]{mathalfa}
\usepackage{xfrac}
\usepackage[unicode,hidelinks,bookmarks=false]{hyperref}
\newcommand{\ie}{i.e.\ }
\newcommand{\cf}{cf.\ }
\newcommand{\resp}{respectively\ }
\newcommand{\Subsection}{\subsection}
\providecommand{\nobreakdash}{\nobreak\hbox{-}}
\setlength{\emergencystretch}{2em}
\multlinegap0pt
\usepackage{amssymb}
\usepackage{bm}
\newtheorem{lem}{Lemma}
\newtheorem{thm}{Theorem}
\newcommand{\C}{\mathbb{C}}
\newcommand{\D}{\mathbb{D}}
\newcommand{\NI}{\noindent}
\newcommand{\T}{\mathbb{T}}
\newcommand{\cln}{\mathcal{N}}
\newcommand{\clr}{\mathcal{R}}
\newcommand{\h}{{H}^2}
\title{Partially isometric truncated and dual truncated Toeplitz operators}
\author{Kritika Babbar, Mo Javed, Amit Maji}
\date{Source: arXiv:2605.21555v1 (CC BY 4.0)}
\begin{document}
\maketitle

\noindent\emph{Source and licence.} Partially isometric truncated and dual truncated Toeplitz operators. Version arXiv:2605.21555v1 (CC BY 4.0), distributed by arXiv under the Creative Commons Attribution 4.0 International licence, http://creativecommons.org/licenses/by/4.0/, as recorded in the arXiv licence metadata for this exact version and shown on the abstract page https://arxiv.org/abs/2605.21555v1. Full licence notice: \texttt{licenses/arxiv-2605.21555-CC-BY-4.0.txt}.\par\smallskip

\noindent\textit{Editorial scope.} Counted unit: the article's thm thm (source0.tex lines 472--474). The article's complete original proof of it is delivered with it (source0.tex lines 477--535, one \texttt{proof} environment of 59 lines). No mathematics is written, repaired or re-derived: the statement and the proof are the article's own lines, in the article's own notation, and no retained line is substituted or re-flowed.  Retained uncounted as the source-local context the proof needs: lines 323--325; lines 328--330; lines 348--355; lines 382--388.  Nothing was omitted from any retained span, so the package carries no omission record.  No retained line contains a citation, so the wrapper carries no bibliography and no bibliography entry.  No retained line contains a figure environment, an image-inclusion command or drawing code, so no image is delivered and the package declares no asset dependency.  Editorial formatting only, in addition: an AMS \texttt{amsart} preamble that restates the article's own declarations and macros used by the retained lines, namely the environments \texttt{lem}, \texttt{thm} on separate counters, together with the standard packages those lines need.  The retained environments print in this document as Lemma 1, Theorem 1 in order of appearance, whereas the article prints the same items with the numbering of its own full document.  The complete native source of the article as distributed in the arXiv e-print for this exact version is bundled unchanged as \texttt{sources/201070/source0.tex}.
\begin{lem}\label{prod}
Let $A_\phi, A_\psi \in \mathscr{T}_{\theta}$ such that $A_\phi A_\psi \in \mathscr{T}_{\theta}$. If one of the operators in the product is in $\mathscr{B}_{\alpha}^{\theta}$ for $\alpha \in \C \cup\{\infty\}$, then either it is a constant multiple of identity or the other is also in $\mathscr{B}_{\alpha}^{\theta}$.
\end{lem}

\begin{lem}\label{typealphachar}
Let $A \in \mathscr{B}_{\alpha}^\theta$ for $\alpha \in \D$. There exists a function $ \psi \in H^\infty$ such that $\|\psi\|_\infty=\|A\|$ and $A=A_{\frac{\psi}{1-\alpha \bar{\theta}}}$. Moreover, if $\phi, \psi \in H^\infty$, then $A_{\frac{\phi}{1-\alpha \bar{\theta}}}A_{\frac{\psi}{1-\alpha \bar{\theta}}}=A_{\frac{\phi\psi}{1-\alpha \bar{\theta}}}$.
\end{lem}

\begin{lem}\label{Lem:A B and D partial isometry equivalences}
Let $\phi\in L^{\infty}(\T)$ with $|\phi|=1$ a.e.\ on $\T$. Then the following are equivalent:
		\begin{enumerate}
			\item $A_{\phi}$ is a partial isometry.
			\item $B_{\phi}$ is a partial isometry.
			\item $D_{\phi}$ is a partial isometry.
		\end{enumerate}
\end{lem}

\begin{thm}\label{Thm:Du partial isometry characterization}
Let $\theta$ be a non-constant inner function and $D_u$ be the dual TTO on $K_{\theta}^{\perp}$ corresponding to the inner function $u$. Then the following are equivalent:
\begin{enumerate}
\item $D_u$ is a partial isometry.
\item $u$ divides $\theta$ or $\theta$ divides $u$.
\end{enumerate}
\end{thm}

\begin{thm}
Let $\theta, u$ be non-constant inner functions such that $\bar{u}\theta$ is not co-analytic. Then $D_{\bar{u} \theta}$ is a partial isometry if and only if $u$ divides $\theta$ or $\theta^2$ divides $u-c$ for some constant c. 
\end{thm}

\begin{proof}
First note that if $D_{\bar{u}\theta}$ is a partial isometry, then $A_{\bar{u}\theta}$ is a partial isometry by Lemma \ref{Lem:A B and D partial isometry equivalences}. Now if $A_{\bar{u}\theta} = 0$, then there exist $h_1, h_2 \in H^2$ such that
$\bar{u}\theta= \theta h_1 +\overline{\theta h_2}$.
Thus
\[
{u}=\overline{h_1}+{\theta^2 h_2}
\]
which implies that $h_1$ is a constant, say $c$. Therefore, $\theta^2 |(u-c)$.
		
Now if $A_{\bar{u} \theta}$ is a non-zero partial isometry, then there exists $0 \neq f \in \cln(A_{\bar{u}\theta})^{\perp}$ such that $\|A_{\bar{u}\theta} \| = \|f \|$.
This implies that $\bar{u}\theta f \in K_\theta$ and hence
\[
T_\theta^* (\bar{u}\theta f)=0.
\]
Now
\[
0=T_\theta^* (\bar{u}\theta f)= T_u^* f, \quad \mbox{and hence} \quad A_{\bar{u}} f =0.
\]
Thus, $\cln (A_{\bar{u}\theta})^\perp \subseteq \cln (A_{\bar{u}})$ which implies 	$\clr(A_u) \subseteq \cln(A_{\bar{u}\theta})$. Consequently, $A_{\bar{u}\theta} A_u=0$. Note that $A_{u}$ is of type $0$. So, by Lemma \ref{prod}, there may have three cases:

\NI		
\textit{Case I:} $A_{u}=0$.

It follows that $u\in\theta \h$ but this  will imply that $\bar{u}\theta$ is co-analytic which is not the case.
\vspace{.1cm}
	
\NI		
\textit{Case II:} $A_{u}= \alpha I$ for some $\alpha \neq 0$.

In this case, we get $A_{\bar{u}\theta}=0$ which is not so.
\vspace{.1cm}

\NI		
\textit{Case III:} $A_{\bar{u}\theta}$ is of type $0$.

It follows by Lemma \ref{typealphachar} that $A_{\bar{u}\theta}=A_\psi$ for some $\psi \in H^\infty$ such that $\|\psi\|_\infty=\|A_{\bar{u}\theta}\|=1$. Thus $A_{\psi}$ is a non-zero partial isometry and hence $\psi$ is inner. Now Theorem \ref{Thm:Du partial isometry characterization} yields either $\theta | \psi$ or $\psi | \theta$. If $\theta | \psi$, then clearly $A_\psi=0$ which is not possible as $A_{\bar{u}\theta}$ is non-zero. Therefore, $\psi | \theta$, i.e., $\theta=\psi w$ for some inner function $w$.
Again $A_{\bar{u}\theta}=A_\psi$ implies that there exist $h_1, h_2 \in \h$ such that
\[
\bar{u}\theta =\psi +\theta h_1 +\overline{\theta h_2}.
\]
Therefore,
\begin{align} \label{eq1}
u &= \bar{\psi}\theta +\overline{h_1} +\theta^2 h_2.
\end{align}
Since $u, \bar{\psi}\theta,$ and $\theta ^2 h_2 \in \h$, $\overline{h_1} \in \h$. 
Hence $h_1$ is constant, say $\beta$. Now (\ref{eq1}) implies 
\begin{align}\label{eq2}
u\psi=\theta + \bar{\beta} \psi +\theta^2 \psi h_2.
\end{align}
Note that 
\[
A_{\psi u}=A_\psi A_u=A_{\bar{u}\theta}A_u=0. 
\]
Then considering the corresponding TTO in (\ref{eq2}) we get $\bar{\beta} A_\psi=0$. Consequently, $\beta=0$. Now from the equation (\ref{eq1}), we get $u=w +\theta^2 h_2$. Again considering the corresponding TTO, we have $A_u=A_w$. Since $w | \theta$, $A_w$ is a partial isometry and hence $A_u$ is so. Also $A_u$ is non-zero, thus $u | \theta$.

Conversely, if $\theta^2|u-c$ for some constant $c$, then $A_{\bar{u}\theta}=0$ and hence a partial isometry. If $u|\theta$, then $\bar{u}\theta$ is inner and also divides 
$\theta$. Therefore, by Theorem \ref{Thm:Du partial isometry characterization},  $D_{\bar{u}\theta}$ is a partial isometry.
This completes the proof.		
\end{proof}

\end{document}
